% ══════════════════════════════════════════════════════════════
%  APPENDIX A: PLAYER INJURIES AND WORKLOAD
% ══════════════════════════════════════════════════════════════

\section{Player Injuries and Workload}
\label{app:injuries}

The welfare concern most frequently raised about the directive is
player health. Longer matches mean more minutes of high-intensity
exertion at the point when fatigue peaks. This appendix reports a
comprehensive analysis using Transfermarkt injury records across
nine leagues and 23,255~player-seasons. Table~\ref{tab:injuries}
presents the results.


\paragraph{Injury counts.}
Total injuries per player-season \textit{decreased} in treated
leagues relative to controls ($\hat{\beta} = -0.090$, SE~$= 0.032$,
$p = 0.005$, pre-treatment mean~$= 1.89$). The decline was driven
entirely by non-muscular injuries---a category dominated by contact
injuries (tackles, collisions) and ligament damage---which fell
significantly ($\hat{\beta} = -0.083$, $p = 0.005$). Muscular
injuries---the fatigue-sensitive category encompassing hamstring
strains, calf tears, and quadriceps injuries---were completely
unchanged ($\hat{\beta} = -0.007$, $p = 0.73$). The share of
injuries classified as muscular was unchanged ($\hat{\beta} = -0.001$,
$p = 0.90$). No positional heterogeneity emerged: muscular injury
rates were unchanged for attackers/midfielders ($\hat{\beta} = -0.011$,
$p = 0.67$) and defenders/goalkeepers
($\hat{\beta} \approx 0$, $p = 0.999$). The null on muscular
injuries is well-powered: the design detects a 10\% effect with
89\% power.

\paragraph{Severity.}
Total days absent was unchanged ($\hat{\beta} = 0.89$, $p = 0.73$,
pre-treatment mean~$= 66.5$ days). Games missed was unchanged
($\hat{\beta} = 0.33$, $p = 0.29$). High-severity injuries were
unchanged ($\hat{\beta} = -0.019$, $p = 0.19$). Injury recurrence
was unchanged ($\hat{\beta} = -0.017$, $p = 0.17$). Per-1,000-hour
injury rates were unchanged for both total ($\hat{\beta} = 0.15$,
$p = 0.93$) and muscular injuries ($\hat{\beta} = 0.24$, $p = 0.51$).

\paragraph{Workload resolution.}
Per-match minutes played rose by $0.76$~minutes ($p < 0.001$),
confirming individual matches are longer. But season-level aggregates
tell the opposite story: total minutes per player-season fell by
70.5~minutes ($p = 0.009$) and total matches fell by 1.01
($p < 0.001$). Average days between matches was unchanged
($p = 0.91$). Clubs responded to longer matches with greater
squad rotation---enabled by the permanent adoption of five
substitutions---keeping cumulative workload roughly constant.

\paragraph{Interpretation.}
The decline in non-muscular injuries likely reflects concurrent
trends---less physical play (documented in the fouls analysis), VAR's
deterrent effect on dangerous challenges, improvements in medical
protocols---rather than a direct effect of the directive. We report it
as context for the muscular-injury null, not as evidence that the
directive improved player safety. The policy-relevant conclusion is
that the directive's benefits (more drama, unchanged competitive
balance) came without measurable cost to player welfare.

% ══════════════════════════════════════════════════════════════
%  TABLE: Player Injuries and Workload
% ══════════════════════════════════════════════════════════════

\begin{table}[t]
\centering
\begin{threeparttable}
\caption{Player Injuries and Workload: Difference-in-Differences}
\label{tab:injuries}
\small
\begin{tabular}{lcccc}
\toprule
& Pre-directive & DiD & & \\
Outcome & mean (treated) & $\hat{\beta}$ & SE & $p$-value \\
\midrule
\multicolumn{5}{l}{\textit{Panel A: Injury counts (per player-season)}} \\[4pt]
Total injuries & 1.886 & $-$0.090 & 0.032 & 0.005 \\
\quad Non-muscular & 1.361 & $-$0.083 & 0.030 & 0.005 \\
\quad Muscular & 0.525 & $-$0.007 & 0.019 & 0.729 \\
Muscular share & 0.270 & $-$0.001 & 0.010 & 0.899 \\[6pt]
\multicolumn{5}{l}{\textit{Panel B: Injury severity (per player-season)}} \\[4pt]
Total days absent & 66.5 & 0.894 & 2.603 & 0.731 \\
Games missed & 8.70 & 0.329 & 0.314 & 0.295 \\
Max severity (days) & 52.5 & $-$1.329 & 2.347 & 0.571 \\
Has recurrence & 0.286 & $-$0.017 & 0.013 & 0.169 \\
High-severity injury & 0.506 & $-$0.019 & 0.014 & 0.190 \\[6pt]
\multicolumn{5}{l}{\textit{Panel C: Per-1,000-hour rates (per player-season)}} \\[4pt]
Total injury rate & 5.77 & 0.146 & 1.557 & 0.925 \\
Muscular injury rate & 1.27 & 0.238 & 0.362 & 0.510 \\[6pt]
\multicolumn{5}{l}{\textit{Panel D: Workload}} \\[4pt]
Minutes per match & 73.7 & 0.764 & 0.145 & $<$0.001 \\
Season total minutes & --- & $-$70.5 & 27.1 & 0.009 \\
Season total matches & --- & $-$1.009 & 0.292 & $<$0.001 \\
Days between matches & --- & $-$0.049 & 0.436 & 0.911 \\
\bottomrule
\end{tabular}
\begin{tablenotes}[flushleft]\footnotesize
\item \textit{Notes:} Panels A--C: $N = 23{,}255$ player-seasons
across nine leagues with Transfermarkt injury data. Panel D:
$N = 830{,}569$ player-match observations (minutes per match);
$N = 21{,}088$ player-seasons (season totals). All specifications
include league and season fixed effects with standard errors
clustered at the league level. Muscular injuries include hamstring,
calf, quadriceps, and adductor injuries. Non-muscular includes
contact injuries, ligament damage, and other trauma. The muscular
injury null is well-powered: design detects a 10\% effect with
89\% power.
\end{tablenotes}
\end{threeparttable}
\end{table}


